\chapter{SM4: Recovering the Master Key from Four Round Keys}
\section{Required subkeys}
If four consecutive final subkeys have been recovered, then
\begin{center}\begin{tabular}{cc}\toprule Key-schedule word&Round key\\\midrule\(K_{32}\)&\(rk_{28}=\hex{428D3654}\)\\\(K_{33}\)&\(rk_{29}=\hex{62293496}\)\\\(K_{34}\)&\(rk_{30}=\hex{01CF72E5}\)\\\(K_{35}\)&\(rk_{31}=\hex{9124A012}\)\\\bottomrule\end{tabular}\end{center}
These are reference subkeys from the official example; \emph{only} \(rk_{31}\) was recovered in the preceding pedagogical DFA walkthrough. the attack walkthrough in Part 2 explains how the current code actually recovers all four subkeys with a selectable-round fault oracle.
\section{Reverse the recurrence}
The forward key schedule is
\[K_{i+4}=K_i\xor T'(K_{i+1}\xor K_{i+2}\xor K_{i+3}\xor CK_i).\]
XOR is self-inverse, so
\[\boxed{K_i=K_{i+4}\xor T'(K_{i+1}\xor K_{i+2}\xor K_{i+3}\xor CK_i)}.\]
Start with \(i=31\), then decrease to \(i=0\). For example,
\[K_{31}=K_{35}\xor T'(K_{32}\xor K_{33}\xor K_{34}\xor CK_{31}),\]
where \(CK_{31}=\hex{646B7279}\). This yields the previous expanded-key word \(K_{31}=\hex{F1780C81}\) in the official test vector. Continue to recover \(K_0,K_1,K_2,K_3\).
\section{Undo the initial constants}
Finally,
\[MK_i=K_i\xor FK_i\qquad(i=0,1,2,3),\]
and concatenate the four master-key words:
\[\boxed{MK=\hex{01234567\allowbreak{}89ABCDEF\allowbreak{}FEDCBA98\allowbreak{}76543210}}.\]
